\documentclass[10pt,a4paper]{amsart}
\usepackage[margin=19mm]{geometry}
\usepackage{amsmath,amssymb,mathtools}
\usepackage[scr=rsfs,cal=euler]{mathalfa}
\usepackage{xfrac}
\usepackage[unicode,hidelinks,bookmarks=false]{hyperref}
\newcommand{\ie}{i.e.\ }
\newcommand{\cf}{cf.\ }
\newcommand{\resp}{respectively\ }
\newcommand{\Subsection}{\subsection}
\providecommand{\nobreakdash}{\nobreak\hbox{-}}
\setlength{\emergencystretch}{2em}
\multlinegap0pt
\usepackage{amssymb}
\usepackage{mathtools}
\usepackage[shortlabels]{enumitem}
\newtheorem{theorem}{Theorem}
\title{Multiscale Localized Inference for Networks\\with a Measured Vertex Coordinate}
\author{Marios Papamichalis, Regina Ruane}
\date{Source: arXiv:2512.18592v2 (CC BY 4.0)}
\begin{document}
\maketitle

\noindent\emph{Source and licence.} Multiscale Localized Inference for Networks\\with a Measured Vertex Coordinate. Version arXiv:2512.18592v2 (CC BY 4.0), distributed by arXiv under the Creative Commons Attribution 4.0 International licence, http://creativecommons.org/licenses/by/4.0/, as recorded in the arXiv licence metadata for this exact version and shown on the abstract page https://arxiv.org/abs/2512.18592v2. Full licence notice: \texttt{licenses/arxiv-2512.18592-CC-BY-4.0.txt}.\par\smallskip

\noindent\textit{Editorial scope.} Counted unit: the article's theorem thm:adaptive-detection-localization (source0.tex lines 5143--5912). The article's complete original proof of it is delivered with it (source0.tex lines 5914--6691, one \texttt{proof} environment of 778 lines). No mathematics is written, repaired or re-derived: the statement and the proof are the article's own lines, in the article's own notation, and no retained line is substituted or re-flowed.  No further source-local context is needed: every cross-reference of every retained line resolves inside the two delivered spans.  Nothing was omitted from any retained span, so the package carries no omission record.  No retained line contains a citation, so the wrapper carries no bibliography and no bibliography entry.  No retained line contains a figure environment, an image-inclusion command or drawing code, so no image is delivered and the package declares no asset dependency.  Editorial formatting only, in addition: an AMS \texttt{amsart} preamble that restates the article's own declarations and macros used by the retained lines, namely the environments \texttt{theorem} on separate counters, together with the standard packages those lines need.  The retained environments print in this document as Theorem 1 in order of appearance, whereas the article prints the same items with the numbering of its own full document.  The complete native source of the article as distributed in the arXiv e-print for this exact version is bundled unchanged as \texttt{sources/191407/source0.tex}.
\begin{theorem}[Global adaptive detection and localization boundary]
\label{thm:adaptive-detection-localization}
Retain the sparse logistic fixed-design model and the notation
\(\alpha_n,\rho_n,B,\mathcal F_{0n},\Delta_n,r_n\) from the preceding
sections, and write
\[
\sigma(x)=\frac{e^x}{1+e^x}.
\]

Let $\mathscr P_n$ denote the statistical model class. Let
\[
\Lambda_n
=
\bigcup_{h\in\mathcal J_n}\Lambda_n(h),
\qquad
m_n=|\Lambda_n|\longrightarrow\infty,
\qquad
\ell_n=\log(2m_n),
\]
be a deterministic, outcome-independent packing of unique symmetric
off-diagonal Haar atoms, with transposed duplicates removed.  We take
\(\mathcal J_n\) to contain only scales \(h\) for which
\(\Lambda_n(h)\ne\varnothing\).  For
\(\lambda\in\Lambda_n\), let \(h_\lambda\) be its scale and put
\[
S_\lambda=\operatorname{supp}(\Psi_\lambda^S).
\]
Here and below, \(|S|\) denotes Lebesgue measure.

Assume
\[
\Psi_\lambda^S\in V_{J_0}^{\perp},
\qquad
\lambda\in\Lambda_n,
\]
and
\begin{equation}
\left\langle
\Psi_\lambda^S,\Psi_\mu^S
\right\rangle_{L^2([0,1]^2)}
=
\mathbf 1\{\lambda=\mu\},
\qquad
\lambda,\mu\in\Lambda_n.
\label{eq:t3-1}
\end{equation}
Assume also that atoms at a common scale have equal support area and,
for some \(\delta_{\rm pack}>0\),
\begin{equation}
\inf_n
\inf_{\substack{\lambda,\mu\in\Lambda_n\\
                 \lambda\ne\mu,\ h_\lambda=h_\mu}}
\frac{|S_\lambda\triangle S_\mu|}{2|S_\lambda|}
\ge\delta_{\rm pack}.
\label{eq:t3-2}
\end{equation}
The infimum over an empty collection is \(+\infty\).

Let
\[
\mathcal B_n
\subset
\left\{
f\in V_{J_0}:
f\text{ is symmetric and mean zero},\
\|f\|_\infty\le B
\right\}
\]
be deterministic, and define
\begin{equation}
\theta_\lambda(f)
=
\left\langle
(I-\Pi_{V_{J_0}})f,\Psi_\lambda^S
\right\rangle_{L^2([0,1]^2)}.
\label{eq:t3-3}
\end{equation}
For \(a>0\) and \(h\in\mathcal J_n\), put
\[
\mathfrak A_n(a,h)
=
\left\{
f_0+\varsigma a h\Psi_\lambda^S:
\begin{array}{l}
f_0\in\mathcal B_n,\quad
\varsigma\in\{-1,1\},\\
\lambda\in\Lambda_n(h),\quad
\|f_0+\varsigma a h\Psi_\lambda^S\|_\infty\le B
\end{array}
\right\}.
\]
By \eqref{eq:t3-1} and the orthogonal decomposition
\(V_{J_0}\oplus V_{J_0}^{\perp}\), every
\(f\in\mathfrak A_n(a,h)\) has a unique representation of this form.
Denote its atom and sign by \(\lambda(f)\) and \(\varsigma(f)\).

Define
\[
\mathfrak F_n^{\rm up}
=
\mathcal B_n
\cup
\bigcup_{h\in\mathcal J_n}
\bigcup_{\substack{a>0:\\
                   \mathfrak A_n(a,h)\ne\varnothing}}
\mathfrak A_n(a,h).
\]

Let \(\mathfrak S_{\rm sym}\) be the measure algebra of symmetric
Lebesgue-measurable subsets of \([0,1]^2\), identified modulo
Lebesgue-null sets and equipped with
\[
d_\triangle(S,T)=|S\triangle T|
\]
and its Borel sigma-field.  This is a standard Borel space.  Define
\[
\mathcal D_n^{\rm dec}
=
\{\partial\}
\sqcup
\bigl(\mathfrak S_{\rm sym}\times\mathcal J_n\bigr),
\]
where \(\partial\) is a cemetery symbol denoting no selected atom and
\(\mathcal J_n\) has the discrete sigma-field.  For
\(D=(S_D,h_D)\ne\partial\), define
\[
d_{\rm loc}(D,\lambda)
=
\min\left\{
1,\
\frac{|S_D\triangle S_\lambda|}{2|S_\lambda|}
+
\mathbf 1\{h_D\ne h_\lambda\}
\right\},
\]
and set
\[
d_{\rm loc}(\partial,\lambda)=1.
\]

Let \(\widehat{\mathcal I}_n\) be the pilot-admitted family defined
by the pilot gate, and put
\[
\widehat\Lambda_n
=
\Lambda_n\cap\widehat{\mathcal I}_n.
\]
Assume that \(\widehat{\mathcal I}_n\), and hence
\(\{\widehat\Lambda_n=\Lambda_n\}\), is
\(\mathcal F_{0n}\)-measurable and that
\begin{equation}
\gamma_n
:=
\sup_{f\in\mathfrak F_n^{\rm up}}
P_f\{\widehat\Lambda_n\ne\Lambda_n\}
=o(1).
\label{eq:t3-4}
\end{equation}
Assume also that
\[
\widehat\Lambda_n=\Lambda_n
\quad\Longrightarrow\quad
\min_{\lambda\in\Lambda_n}\widehat s_\lambda>0
\quad\text{almost surely}.
\]
If standard-error positivity is instead imposed as a separate gate,
that gate must be included in the admission event in
\eqref{eq:t3-4}.

Suppose that there is an \(\mathcal F_{0n}\)-measurable event
\(\mathcal G_n\subseteq\{\Delta_n\le r_n\}\) such that
\begin{equation}
\sup_{f\in\mathfrak F_n^{\rm up}}
P_f(\mathcal G_n^c)
\le\eta_n+o(1),
\qquad
\eta_n\to0.
\label{eq:t3-5}
\end{equation}
Define
\[
\widetilde{\mathcal G}_n
=
\mathcal G_n
\cap
\{\widehat\Lambda_n=\Lambda_n\}.
\]
Then
\begin{equation}
\sup_{f\in\mathfrak F_n^{\rm up}}
P_f(\widetilde{\mathcal G}_n^c)
\le
\eta_n+\gamma_n+o(1).
\label{eq:t3-6}
\end{equation}

For a random array \(R_{n,f}\), an
\(\mathcal F_{0n}\)-measurable event \(H_n\), and a deterministic
sequence \(q_n>0\), write
\[
R_{n,f}=o_{\rm ucp}(q_n)
\quad
\text{on \(H_n\), uniformly over \(\mathfrak F_n^{\rm up}\),}
\]
if, for every \(\epsilon,\delta>0\),
\[
\sup_{f\in\mathfrak F_n^{\rm up}}
P_f\left[
H_n
\cap
\left\{
P_f\left(
|R_{n,f}|>\epsilon q_n
\mid\mathcal F_{0n}
\right)>\delta
\right\}
\right]
\longrightarrow0.
\]

For \(f\in\mathfrak F_n^{\rm up}\), let
\(\theta_{\lambda,n}^{D}(f)\) be the frozen finite-design target and
put
\[
b_{\lambda,n}(f)
=
\theta_{\lambda,n}^{D}(f)-\theta_\lambda(f).
\]
Let \(s_{\lambda,f}^2\) be the conditional variance of the frozen
first-order influence sum and put
\[
I_{\lambda,f}=s_{\lambda,f}^{-2}.
\]
Assume that
\(\theta_{\lambda,n}^{D}(f)\), \(b_{\lambda,n}(f)\),
\(s_{\lambda,f}\), and \(I_{\lambda,f}\) are
\(\mathcal F_{0n}\)-measurable.  On
\(\widetilde{\mathcal G}_n\), assume uniformly over
\(f\in\mathfrak F_n^{\rm up}\) and
\(\lambda\in\Lambda_n\) that
\begin{equation}
0<\underline\iota n^2\rho_n
\le
I_{\lambda,f}
\le
\overline\iota n^2\rho_n
<\infty
\label{eq:t3-7}
\end{equation}
for constants
\(0<\underline\iota\le\overline\iota<\infty\).

Write
\[
p_{ij}(f)
=
E_f(A_{ij}\mid\mathcal F_{0n}).
\]
For each \(f\in\mathfrak F_n^{\rm up}\), assume that there are
\(\mathcal F_{0n}\)-measurable loadings
\(\gamma_{\lambda,ij}(f)\) such that
\[
\zeta_{\lambda,ij}(f)
=
s_{\lambda,f}^{-1}
\gamma_{\lambda,ij}(f)
\{A_{ij}-p_{ij}(f)\}.
\]
For every fixed \(\lambda\), conditional on \(\mathcal F_{0n}\),
assume that
\(\{\zeta_{\lambda,ij}(f):i<j\}\) are independent and, uniformly over
\(f\in\mathfrak F_n^{\rm up}\), satisfy
\begin{equation}
E_f\{\zeta_{\lambda,ij}(f)\mid\mathcal F_{0n}\}=0,
\qquad
\sum_{i<j}
E_f\{\zeta_{\lambda,ij}^2(f)\mid\mathcal F_{0n}\}=1.
\label{eq:t3-8}
\end{equation}
Dependence across different \(\lambda\)'s is unrestricted.  Assume,
for a deterministic sequence \(L_n\), that, for every
\(f\in\mathfrak F_n^{\rm up}\),
\begin{equation}
\mathbf 1_{\widetilde{\mathcal G}_n}
\max_{\lambda,i<j}
|\zeta_{\lambda,ij}(f)|
\le L_n
\quad P_f\text{-almost surely},
\qquad
L_n^2\log^7(m_n n)\to0.
\label{eq:t3-9}
\end{equation}

Define
\[
R_{n,f}^{\rm lin}
=
\max_{\lambda\in\Lambda_n}
\left|
\frac{
\widehat\theta_\lambda-\theta_{\lambda,n}^{D}(f)
}{
s_{\lambda,f}
}
-
\sum_{i<j}\zeta_{\lambda,ij}(f)
\right|.
\]
Assume
\begin{equation}
R_{n,f}^{\rm lin}
=
o_{\rm ucp}(\ell_n^{-1/2})
\quad
\text{on }\widetilde{\mathcal G}_n.
\label{eq:t3-10}
\end{equation}
Also define
\[
\delta_{n,f}
=
\max_{\lambda\in\Lambda_n}
\left|
\frac{\widehat s_\lambda}{s_{\lambda,f}}-1
\right|,
\]
and assume
\begin{equation}
\delta_{n,f}
=
o_{\rm ucp}(\ell_n^{-1})
\quad
\text{on }\widetilde{\mathcal G}_n.
\label{eq:t3-11}
\end{equation}

Assume, for every \(\epsilon>0\),
\begin{equation}
\sup_{f\in\mathfrak F_n^{\rm up}}
P_f\left[
\widetilde{\mathcal G}_n
\cap
\left\{
\sqrt{\ell_n}
\max_{\lambda\in\Lambda_n}
\frac{|b_{\lambda,n}(f)|}{s_{\lambda,f}}
>\epsilon
\right\}
\right]
\longrightarrow0.
\label{eq:t3-12}
\end{equation}
Here \(b_{\lambda,n}(f)\) contains only frozen,
\(\mathcal F_{0n}\)-measurable target discrepancies.  Stochastic
Taylor and nonlinear-estimation remainders are included in
\eqref{eq:t3-10}, not in \(b_{\lambda,n}(f)\).

A sufficient primitive verification of \eqref{eq:t3-12} is as
follows.  Decompose
\[
b_{\lambda,n}(f)
=
b_{\lambda,n}^{\rm nonord}(f)
+
b_{\lambda,n}^{\rm order}(f).
\]
Assume that the aggregate nonordering term satisfies the analogue of
\eqref{eq:t3-12} and that, for every
\(f\in\mathfrak F_n^{\rm up}\), on
\(\widetilde{\mathcal G}_n\),
\[
\max_{\lambda\in\Lambda_n}
|b_{\lambda,n}^{\rm order}(f)|
\le
C_{\rm ord}Br_n.
\]
Then
\begin{equation}
Br_n\,n\sqrt{\rho_n\ell_n}
\longrightarrow0
\label{eq:t3-13}
\end{equation}
is sufficient for the ordering term to satisfy the analogue of
\eqref{eq:t3-12}.  Indeed, by \eqref{eq:t3-7},
\[
\sqrt{\ell_n}
\max_\lambda
\frac{|b_{\lambda,n}^{\rm order}(f)|}{s_{\lambda,f}}
\le
C_{\rm ord}\sqrt{\overline\iota}\,
Br_n n\sqrt{\rho_n\ell_n}.
\]
For any sequences \(a_n,h_n\) satisfying the upper signal condition
\eqref{eq:t3-17} below,
\[
\frac{Br_n}{a_n h_n}
\le
\frac{
Br_n n\sqrt{\rho_n\ell_n}
}{
\sqrt{C_{\rm up}}\,\ell_n
}
\longrightarrow0.
\]
Thus \(r_n/h_n\le c_r\) alone is insufficient when \(a_n\to0\).

Let \(\mathcal X_n\) denote the full observed data.  Conditionally on
\(\mathcal X_n\), assume that
\[
Z^{*(b)}
=
\bigl(
Z_\lambda^{*(b)}:
\lambda\in\Lambda_n
\bigr),
\qquad b=1,\ldots,B_n^*,
\]
are independent centered jointly Gaussian vectors satisfying
\begin{equation}
\operatorname{Var}^*(Z_\lambda^{*(b)})=1,
\qquad
\lambda\in\Lambda_n.
\label{eq:t3-14}
\end{equation}
Dependence among coordinates within the same multiplier vector is
unrestricted.  Let \(P_{\mathcal X_n}^*\) be the conditional law of
one multiplier vector given \(\mathcal X_n\), and let
\(P_{\mathcal X_n}^{*,B}\) be the conditional product law of all
\(B_n^*\) multiplier vectors.  Assume \(B_n^*\to\infty\), and define
\[
T^{*(b)}
=
\max_{\lambda\in\Lambda_n}|Z_\lambda^{*(b)}|,
\qquad
F_n^*(t)
=
P_{\mathcal X_n}^*(T^{*(1)}\le t),
\]
\[
q_{1-u,n}^*
=
\inf\{t:F_n^*(t)\ge1-u\},
\]
and
\[
\widehat F_{B_n^*}^*(t)
=
\frac1{B_n^*}
\sum_{b=1}^{B_n^*}
\mathbf 1\{T^{*(b)}\le t\},
\]
\[
c_{1-\alpha}^*
=
\inf\left\{
t:
\widehat F_{B_n^*}^*(t)\ge1-\alpha
\right\},
\qquad
\alpha\in(0,1).
\]
Let \(P_f^\dagger\) and \(E_f^\dagger\) denote the joint law of the
data and the independent Gaussian seeds used for the multiplier
draws.

On \(\{\widehat\Lambda_n=\Lambda_n\}\), define
\[
Z_\lambda
=
\frac{\widehat\theta_\lambda}{\widehat s_\lambda},
\qquad
T_n
=
\max_{\lambda\in\Lambda_n}|Z_\lambda|.
\]
When admission fails, set
\[
Z_\lambda=0
\quad(\lambda\in\Lambda_n),
\qquad
T_n=0.
\]
Define
\[
\phi_n
=
\mathbf 1\{\widehat\Lambda_n=\Lambda_n\}
\mathbf 1\{T_n>c_{1-\alpha}^*\}.
\]
With deterministic tie-breaking, put
\[
\widehat\lambda
=
\begin{cases}
\displaystyle
\arg\max_{\lambda\in\Lambda_n}|Z_\lambda|,
&\phi_n=1,\\[2mm]
\varnothing,&\phi_n=0,
\end{cases}
\]
and
\[
(\widehat\varepsilon,\widehat D)
=
\begin{cases}
\left(
\operatorname{sign}(\widehat\theta_{\widehat\lambda}),
(S_{\widehat\lambda},h_{\widehat\lambda})
\right),
&\widehat\lambda\ne\varnothing,\\[1mm]
(0,\partial),
&\widehat\lambda=\varnothing.
\end{cases}
\]

For \(f_0\in\mathcal B_n\), define
\[
F_{n,f_0}^{0}(t)
=
P_{f_0}(T_n\le t\mid\mathcal F_{0n}),
\qquad
\Delta_{n,f_0}^{\rm cal}
=
\sup_{t\in\mathbb R}
|F_{n,f_0}^{0}(t)-F_n^*(t)|.
\]
Assume the following uniform calibration condition, as established
by the preceding multiplier theorem: there is a deterministic
\(\kappa_n\downarrow0\) such that
\begin{equation}
\sup_{f_0\in\mathcal B_n}
P_{f_0}\left[
\widetilde{\mathcal G}_n
\cap
\{\Delta_{n,f_0}^{\rm cal}>\kappa_n\}
\right]
\longrightarrow0.
\label{eq:t3-15}
\end{equation}
The calibration condition is required for the actual statistic
\(T_n\), the deterministic family \(\Lambda_n\), and the same
loadings and studentizers as the implemented scan.  Then
\begin{equation}
\sup_{f_0\in\mathcal B_n}
P_{f_0}^\dagger(\phi_n=1)
\le
\alpha+\eta_n+o(1).
\label{eq:t3-16}
\end{equation}
This is a global-null size statement.  Strong familywise error
control requires the corresponding subset-null calibration result.

There exists \(C_{\rm up}<\infty\) such that the following holds.
For every deterministic sequence
\[
h_n\in\mathcal J_n,
\qquad
a_n>0,
\qquad
\mathfrak A_n(a_n,h_n)\ne\varnothing,
\]
satisfying
\begin{equation}
a_n^2n^2\rho_nh_n^2
\ge
C_{\rm up}\ell_n,
\label{eq:t3-17}
\end{equation}
define
\[
\mathcal E_{n,f}^{\rm rec}
=
\left\{
\phi_n=0
\ \text{or}\
\widehat\lambda\ne\lambda(f)
\ \text{or}\
\widehat\varepsilon\ne\varsigma(f)
\right\}.
\]
Then, for every \(\delta>0\),
\begin{equation}
\sup_{f\in\mathfrak A_n(a_n,h_n)}
P_f\left[
\widetilde{\mathcal G}_n
\cap
\left\{
P_f^\dagger
\left(
\mathcal E_{n,f}^{\rm rec}
\mid\mathcal F_{0n}
\right)
>\delta
\right\}
\right]
\longrightarrow0.
\label{eq:t3-18}
\end{equation}
Consequently,
\begin{equation}
\sup_{f\in\mathfrak A_n(a_n,h_n)}
E_f^\dagger
d_{\rm loc}\{\widehat D,\lambda(f)\}
\le
\eta_n+\gamma_n+o(1)
\longrightarrow0.
\label{eq:t3-19}
\end{equation}

For the lower bound, all probabilities, expectations, and risks below
are conditional on a deterministic design
\((U_1,\ldots,U_n)\) satisfying the following conditions.  Suppose
that \(\mathcal B_n\) contains a symmetric, mean-zero interior
element \(f_0^\circ\) and constants \(\delta_B,c_p,C_p>0\) such that
\begin{equation}
\|f_0^\circ\|_\infty\le B-\delta_B,
\qquad
c_p\rho_n\le p_{0,ij}\le C_p\rho_n,
\qquad
1-p_{0,ij}\ge c_p,
\label{eq:t3-20}
\end{equation}
where
\[
p_{0,ij}
=
\sigma\{\alpha_n+f_0^\circ(U_i,U_j)\}.
\]

Let
\[
\mathcal K_n
=
\{(i,j):1\le i<j\le n\},
\qquad
v_{\lambda,ij}
=
h_\lambda\Psi_\lambda^S(U_i,U_j).
\]
Assume
\begin{equation}
\|v_\lambda\|_\infty\le C_\Psi,
\qquad
c_Xn^2h_\lambda^2
\le
\sum_{(i,j)\in\mathcal K_n}v_{\lambda,ij}^2
\le
C_Xn^2h_\lambda^2
\label{eq:t3-21}
\end{equation}
uniformly over \(\lambda\in\Lambda_n\).  The two-sided energy
bound records membership in the same regular-design class; only its
upper bound is used below.

Let
\[
\mathbf a_n=(a_h:h\in\mathcal J_n),
\qquad
a_h>0,
\]
and assume
\begin{equation}
\max_{h\in\mathcal J_n}
a_h
\max_{\lambda\in\Lambda_n(h)}
\|h\Psi_\lambda^S\|_\infty
\le\delta_B.
\label{eq:t3-22}
\end{equation}
For
\[
\mathcal D_\lambda
=
\{(i,j)\in\mathcal K_n:v_{\lambda,ij}\ne0\},
\]
define
\[
N_{{\rm ov},n}
=
\#\left\{
(\lambda,\mu)\in\Lambda_n^2:
\mathcal D_\lambda\cap\mathcal D_\mu\ne\varnothing
\right\},
\]
including every label-diagonal pair \((\lambda,\lambda)\).  Assume
\begin{equation}
N_{{\rm ov},n}
\le m_nd_n,
\qquad
d_n\ge1,
\qquad
\frac{\log d_n}{\log(2m_n)}
\longrightarrow0.
\label{eq:t3-23}
\end{equation}

Let \(\mathcal E_n^{\rm orc}\) be the experiment observing the
deterministic design, \(\alpha_n\), the fixed baseline
\(f_0^\circ\), and all outcomes
\(\{A_{ij}:(i,j)\in\mathcal K_n\}\).  Assume that, after restricting
the baseline to \(f_0^\circ\), the original experiment is obtained
from \(\mathcal E_n^{\rm orc}\) through a Markov kernel that does not
depend on \(\lambda\), the sign, or the amplitude.  This includes
outcome-independent missingness conditional on the fixed design.

Define
\[
\Theta_{1,n}(\mathbf a_n)
=
\left\{
\vartheta=(f_0,\varsigma,h,\lambda):
\begin{array}{l}
f_0\in\mathcal B_n,\quad
\varsigma\in\{-1,1\},\quad
h\in\mathcal J_n,\\
\lambda\in\Lambda_n(h),\quad
\|f_0+\varsigma a_h h\Psi_\lambda^S\|_\infty\le B
\end{array}
\right\},
\]
and let
\[
f_\vartheta
=
f_0+\varsigma a_h h\Psi_\lambda^S,
\qquad
\lambda(\vartheta)=\lambda.
\]
Let \(P_\vartheta\) and \(E_\vartheta\) denote the conditional law and
expectation in the original experiment.  Define
\begin{align}
R_n^{\rm det}
&=
\inf_\varphi
\left[
\sup_{f_0\in\mathcal B_n}P_{f_0}(\varphi=1)
+
\sup_{\vartheta\in\Theta_{1,n}(\mathbf a_n)}
P_\vartheta(\varphi=0)
\right],
\label{eq:t3-24}\\
R_n^{\rm loc}
&=
\inf_{\widetilde D}
\sup_{\vartheta\in\Theta_{1,n}(\mathbf a_n)}
E_\vartheta
d_{\rm loc}\{\widetilde D,\lambda(\vartheta)\},
\label{eq:t3-25}
\end{align}
where randomized tests are allowed and the second infimum is over
all possibly randomized measurable decisions taking values in
\(\mathcal D_n^{\rm dec}\).

There exists \(c_{\rm low}>0\) such that, if
\begin{equation}
\max_{h\in\mathcal J_n}
a_h^2n^2\rho_nh^2
\le
c_{\rm low}\ell_n,
\label{eq:t3-26}
\end{equation}
then
\begin{equation}
R_n^{\rm det}\longrightarrow1,
\qquad
\liminf_{n\to\infty}R_n^{\rm loc}>0.
\label{eq:t3-27}
\end{equation}
These are global-union minimax conclusions.
\end{theorem}

\begin{proof}
We first prove the upper assertions.  Fix
\[
f
=
f_0+\varsigma_0a_n h_n\Psi_{\lambda_0}^S
\in\mathfrak A_n(a_n,h_n),
\]
where
\[
\lambda_0=\lambda(f),
\qquad
\varsigma_0=\varsigma(f).
\]
By \eqref{eq:t3-1}--\eqref{eq:t3-3},
\begin{equation}
\theta_{\lambda_0}(f)
=
\varsigma_0a_n h_n,
\qquad
\theta_\mu(f)=0
\quad(\mu\ne\lambda_0).
\label{eq:t3-28}
\end{equation}

For notational brevity, write
\[
G_\lambda
=
\sum_{i<j}\zeta_{\lambda,ij}(f)
\]
and put
\[
r_{n,f}
=
\max_{\lambda\in\Lambda_n}
\left|
\frac{\widehat\theta_\lambda-\theta_\lambda(f)}
{s_{\lambda,f}}
-
G_\lambda
\right|.
\]
By the definition of \(b_{\lambda,n}(f)\),
\[
r_{n,f}
\le
R_{n,f}^{\rm lin}
+
\max_{\lambda\in\Lambda_n}
\frac{|b_{\lambda,n}(f)|}{s_{\lambda,f}}.
\]
The second summand is \(\mathcal F_{0n}\)-measurable.  Therefore its
conditional exceedance probability is either zero or one, and
\eqref{eq:t3-10}--\eqref{eq:t3-12} imply
\begin{equation}
r_{n,f}
=
o_{\rm ucp}(\ell_n^{-1/2})
\quad
\text{on }\widetilde{\mathcal G}_n.
\label{eq:t3-29}
\end{equation}

For every fixed \(\lambda\), conditional Bernstein's inequality gives
\[
P_f\left(
|G_\lambda|>x
\mid\mathcal F_{0n}
\right)
\le
2\exp\left\{
-\frac{x^2}{2(1+L_nx/3)}
\right\}
\]
on \(\widetilde{\mathcal G}_n\).  Hence, for every
\(K_0>\sqrt2\),
\begin{align}
&P_f\left(
\max_{\lambda\in\Lambda_n}|G_\lambda|
>
K_0\sqrt{\ell_n}
\mid\mathcal F_{0n}
\right)
\notag\\
&\qquad\le
2m_n
\exp\left\{
-\frac{K_0^2\ell_n}
{2\{1+K_0L_n\sqrt{\ell_n}/3\}}
\right\}
=o(1),
\label{eq:t3-30}
\end{align}
uniformly on \(\widetilde{\mathcal G}_n\).  Here
\(L_n\sqrt{\ell_n}\to0\) follows from \eqref{eq:t3-9}.  On the
complement of the event in \eqref{eq:t3-30},
\begin{equation}
\max_{\lambda\in\Lambda_n}
\frac{
|\widehat\theta_\lambda-\theta_\lambda(f)|
}{
s_{\lambda,f}
}
\le
K_0\sqrt{\ell_n}+r_{n,f}.
\label{eq:t3-31}
\end{equation}

For one ideal multiplier draw, conditional Gaussianity gives
\[
P_{\mathcal X_n}^*(T^{*(1)}>x)
\le
2m_ne^{-x^2/2}.
\]
Let
\[
x_{\alpha,n}
=
\sqrt{2\log(4m_n/\alpha)}.
\]
Then
\[
P_{\mathcal X_n}^*
(T^{*(1)}>x_{\alpha,n})
\le\alpha/2.
\]
Conditional on \(\mathcal X_n\), the number of the \(B_n^*\) draws
exceeding \(x_{\alpha,n}\) is binomial with success probability at
most \(\alpha/2\).  Thus, for some \(c_\alpha>0\),
\begin{equation}
P_{\mathcal X_n}^{*,B}
\{c_{1-\alpha}^*>x_{\alpha,n}\}
\le
e^{-c_\alpha B_n^*}.
\label{eq:t3-32}
\end{equation}
Since
\[
x_{\alpha,n}/\sqrt{\ell_n}\longrightarrow\sqrt2,
\]
every fixed \(K_c>\sqrt2\) satisfies, for all sufficiently large \(n\),
\begin{align}
&P_f^\dagger\left(
c_{1-\alpha}^*>K_c\sqrt{\ell_n}
\mid\mathcal F_{0n}
\right)
\notag\\
&\quad=
E_f\left[
P_{\mathcal X_n}^{*,B}
\{c_{1-\alpha}^*>K_c\sqrt{\ell_n}\}
\mid\mathcal F_{0n}
\right]
\notag\\
&\quad\le
E_f\left[
P_{\mathcal X_n}^{*,B}
\{c_{1-\alpha}^*>x_{\alpha,n}\}
\mid\mathcal F_{0n}
\right]
\le
e^{-c_\alpha B_n^*}.
\label{eq:t3-33}
\end{align}

The standardized noncentrality is
\[
\nu_n
=
\frac{a_n h_n}{s_{\lambda_0,f}}
=
a_n h_n\sqrt{I_{\lambda_0,f}}.
\]
By \eqref{eq:t3-7} and \eqref{eq:t3-17},
\[
\nu_n^2
\ge
\underline\iota C_{\rm up}\ell_n.
\]
Choose
\[
K_s>\max\{2K_0,K_0+K_c\},
\qquad
C_{\rm up}>
\frac{K_s^2}{\underline\iota}.
\]
Then
\begin{equation}
\nu_n
\ge
K_s\sqrt{\ell_n}.
\label{eq:t3-34}
\end{equation}

On
\[
\left\{
\max_\lambda|G_\lambda|
\le K_0\sqrt{\ell_n}
\right\}
\cap
\{\delta_{n,f}<1\},
\]
we have the exact bounds
\begin{align}
|Z_{\lambda_0}|
&\ge
\frac{
\nu_n-K_0\sqrt{\ell_n}-r_{n,f}
}{
1+\delta_{n,f}
},
\label{eq:t3-35}\\
\max_{\mu\ne\lambda_0}|Z_\mu|
&\le
\frac{
K_0\sqrt{\ell_n}+r_{n,f}
}{
1-\delta_{n,f}
}.
\label{eq:t3-36}
\end{align}

The preceding choice implies both
\(K_s-K_0>K_c\) and \(K_s-K_0>K_0\).
By continuity at zero, there is a constant
\(\bar\delta\in(0,1)\) such that
\[
\frac{K_s-K_0}{1+\bar\delta}
>
\max\left\{
K_c,\,
\frac{K_0}{1-\bar\delta}
\right\}.
\]
For all sufficiently large \(n\), the conditional failure probability
of simultaneous rejection, exact atom recovery, and sign recovery is
bounded by
\begin{align*}
&
P_f\left(
\max_\lambda|G_\lambda|>K_0\sqrt{\ell_n}
\mid\mathcal F_{0n}
\right)
+
P_f(r_{n,f}>1\mid\mathcal F_{0n})\\
&\qquad
+
P_f(\delta_{n,f}>\bar\delta\mid\mathcal F_{0n})
+
e^{-c_\alpha B_n^*}.
\end{align*}
By \eqref{eq:t3-29}, \eqref{eq:t3-30}, \eqref{eq:t3-11}, and
\eqref{eq:t3-33}, this bound converges to zero in the uniform
conditional sense defined before \eqref{eq:t3-7}.  On its complement,
\[
|Z_{\lambda_0}|
>
\max\left\{
c_{1-\alpha}^*,\
\max_{\mu\ne\lambda_0}|Z_\mu|
\right\},
\]
so
\[
\phi_n=1,
\qquad
\widehat\lambda=\lambda_0.
\]
Moreover,
\[
\frac{
\varsigma_0\widehat\theta_{\lambda_0}
}{
s_{\lambda_0,f}
}
\ge
\nu_n-K_0\sqrt{\ell_n}-r_{n,f}>0,
\]
so \(\widehat\varepsilon=\varsigma_0\).  This proves
\eqref{eq:t3-18}.

We next prove \eqref{eq:t3-16}.  Put
\[
\epsilon_n^*=(B_n^*)^{-1/4}.
\]
For all sufficiently large \(n\),
\[
\alpha+\epsilon_n^*+\kappa_n<1.
\]
Define
\[
q_{p,n}^{0}(f_0)
=
\inf\{t:F_{n,f_0}^{0}(t)\ge p\}.
\]
On
\(\{\Delta_{n,f_0}^{\rm cal}\le\kappa_n\}\),
\[
q_{1-\alpha-\epsilon_n^*,n}^*
\ge
q_{1-\alpha-\epsilon_n^*-\kappa_n,n}^{0}(f_0).
\]
Indeed, if
\(t<q_{1-\alpha-\epsilon_n^*-\kappa_n,n}^{0}(f_0)\), then
\[
F_n^*(t)
\le
F_{n,f_0}^{0}(t)+\kappa_n
<
1-\alpha-\epsilon_n^*.
\]
Consequently,
\begin{align*}
&
\left\{
\widetilde{\mathcal G}_n,\
T_n>q_{1-\alpha-\epsilon_n^*,n}^*,\
\Delta_{n,f_0}^{\rm cal}\le\kappa_n
\right\}\\
&\qquad\subseteq
\left\{
\widetilde{\mathcal G}_n,\
T_n>
q_{1-\alpha-\epsilon_n^*-\kappa_n,n}^{0}(f_0)
\right\}.
\end{align*}
Since \(\widetilde{\mathcal G}_n\) and the conditional null quantile
are \(\mathcal F_{0n}\)-measurable, \eqref{eq:t3-15} gives
\[
\sup_{f_0\in\mathcal B_n}
P_{f_0}\left[
\widetilde{\mathcal G}_n
\cap
\left\{
T_n>
q_{1-\alpha-\epsilon_n^*,n}^*
\right\}
\right]
\le
\alpha+\epsilon_n^*+\kappa_n+o(1).
\]

Conditional on \(\mathcal X_n\), DKW gives
\[
P_{\mathcal X_n}^{*,B}\left(
\sup_t
|\widehat F_{B_n^*}^*(t)-F_n^*(t)|
>
\epsilon_n^*
\right)
\le
2e^{-2\sqrt{B_n^*}}.
\]
On the complementary event,
\[
c_{1-\alpha}^*
\ge
q_{1-\alpha-\epsilon_n^*,n}^*.
\]
Therefore
\[
\sup_{f_0\in\mathcal B_n}
P_{f_0}^\dagger\left[
\widetilde{\mathcal G}_n
\cap
\{T_n>c_{1-\alpha}^*\}
\right]
\le
\alpha+o(1).
\]
Admission failure forces \(\phi_n=0\), so
\[
P_{f_0}^\dagger(\phi_n=1)
\le
P_{f_0}^\dagger\left[
\widetilde{\mathcal G}_n
\cap
\{T_n>c_{1-\alpha}^*\}
\right]
+
P_{f_0}(\mathcal G_n^c).
\]
Equation \eqref{eq:t3-5} proves \eqref{eq:t3-16}.

Finally, put
\[
q_{n,f}
=
P_f^\dagger
(\mathcal E_{n,f}^{\rm rec}\mid\mathcal F_{0n}).
\]
For every fixed \(\delta>0\),
\begin{align*}
E_f^\dagger
\left[
\mathbf 1_{\widetilde{\mathcal G}_n}
\mathbf 1_{\mathcal E_{n,f}^{\rm rec}}
\right]
&=
E_f\left[
\mathbf 1_{\widetilde{\mathcal G}_n}q_{n,f}
\right]\\
&\le
\delta
+
P_f\left[
\widetilde{\mathcal G}_n
\cap
\{q_{n,f}>\delta\}
\right].
\end{align*}
Taking the supremum, then the \(\limsup\), and then letting
\(\delta\downarrow0\), \eqref{eq:t3-18} gives
\[
\sup_{f\in\mathfrak A_n(a_n,h_n)}
P_f^\dagger(\mathcal E_{n,f}^{\rm rec})
\le
\eta_n+\gamma_n+o(1).
\]
Since the loss is bounded by one and vanishes under exact recovery,
\eqref{eq:t3-19} follows.

We now prove the lower assertions.  Restrict the null and alternative
classes to the baseline \(f_0^\circ\), and put
\[
P_0:=P_{f_0^\circ}^{\rm orc}.
\]
Every mixture element below belongs to
\(\Theta_{1,n}(\mathbf a_n)\), because
\[
\left\|
f_0^\circ+\varsigma a_h h\Psi_\lambda^S
\right\|_\infty
\le
(B-\delta_B)+\delta_B
=
B.
\]
It remains symmetric and mean zero because \(f_0^\circ\) and the Haar
atoms have these properties.

Let \(R_{n,\circ}^{\rm det,orig}\) and
\(R_{n,\circ}^{\rm loc,orig}\) denote the original-experiment risks
after fixing \(f_0=f_0^\circ\), and let
\(R_{n,\circ}^{\rm det,orc}\) and
\(R_{n,\circ}^{\rm loc,orc}\) be the corresponding oracle risks.
Parameter restriction and Blackwell dominance give
\[
R_n^{\rm det}
\ge
R_{n,\circ}^{\rm det,orig}
\ge
R_{n,\circ}^{\rm det,orc},
\qquad
R_n^{\rm loc}
\ge
R_{n,\circ}^{\rm loc,orig}
\ge
R_{n,\circ}^{\rm loc,orc}.
\]

For \(\lambda\in\Lambda_n\) and
\(\varsigma\in\{-1,1\}\), define
\[
t_{\lambda,\varsigma,ij}
=
\varsigma a_{h_\lambda}v_{\lambda,ij},
\qquad
p_{\lambda,\varsigma,ij}
=
\sigma\left\{
\operatorname{logit}(p_{0,ij})
+
t_{\lambda,\varsigma,ij}
\right\}.
\]
Let \(P_{\lambda,\varsigma}\) be the corresponding product-Bernoulli
oracle law and put
\[
L_{\lambda,\varsigma}
=
\frac{dP_{\lambda,\varsigma}}{dP_0},
\qquad
\Delta_{\lambda,\varsigma,ij}
=
p_{\lambda,\varsigma,ij}-p_{0,ij}.
\]
Direct calculation gives
\begin{equation}
E_0
(L_{\lambda,\varsigma}L_{\mu,\varsigma'})
=
\prod_{(i,j)\in\mathcal K_n}
\left[
1+
\frac{
\Delta_{\lambda,\varsigma,ij}
\Delta_{\mu,\varsigma',ij}
}{
p_{0,ij}(1-p_{0,ij})
}
\right].
\label{eq:t3-37}
\end{equation}

For every fixed \(M<\infty\), there is \(C_M<\infty\) such that
\begin{equation}
|\sigma(\eta+u)-\sigma(\eta)|
\le
C_M\sigma(\eta)\{1-\sigma(\eta)\}|u|,
\qquad |u|\le M.
\label{eq:t3-38}
\end{equation}
This follows because
\[
\sup_{\eta\in\mathbb R}
\sup_{|r|\le M}
\frac{\sigma'(\eta+r)}{\sigma'(\eta)}
<\infty.
\]
By \eqref{eq:t3-20}--\eqref{eq:t3-22},
\(|t_{\lambda,\varsigma,ij}|\le\delta_B\).  Hence
\begin{align}
\log E_0L_{\lambda,\varsigma}^2
&\le
C a_{h_\lambda}^2n^2\rho_nh_\lambda^2,
\label{eq:t3-39}\\
\operatorname{KL}(P_{\lambda,\varsigma},P_0)
&\le
C a_{h_\lambda}^2n^2\rho_nh_\lambda^2.
\label{eq:t3-40}
\end{align}
Indeed,
\[
\log E_0L_{\lambda,\varsigma}^2
\le
\sum_{(i,j)\in\mathcal K_n}
\frac{\Delta_{\lambda,\varsigma,ij}^2}
{p_{0,ij}(1-p_{0,ij})},
\]
and the same sum bounds the product-Bernoulli KL divergence.  By
\eqref{eq:t3-38},
\[
\frac{\Delta_{\lambda,\varsigma,ij}^2}
{p_{0,ij}(1-p_{0,ij})}
\le
C\rho_n a_{h_\lambda}^2v_{\lambda,ij}^2.
\]
Summation and \eqref{eq:t3-21} prove
\eqref{eq:t3-39}--\eqref{eq:t3-40}.

If
\(\mathcal D_\lambda\cap\mathcal D_\mu=\varnothing\), every factor in
\eqref{eq:t3-37} equals one, so
\begin{equation}
E_0
(L_{\lambda,\varsigma}L_{\mu,\varsigma'})
=1.
\label{eq:t3-41}
\end{equation}
For an overlapping pair, put
\[
x_{ij}
=
\frac{
\Delta_{\lambda,\varsigma,ij}
\Delta_{\mu,\varsigma',ij}
}{
p_{0,ij}(1-p_{0,ij})
}.
\]
Each \(1+x_{ij}>0\), so
\(\log(1+x_{ij})\le x_{ij}\le|x_{ij}|\).  Thus
\begin{align}
\log E_0
(L_{\lambda,\varsigma}L_{\mu,\varsigma'})
&\le
C\rho_n a_{h_\lambda}a_{h_\mu}
\sum_{(i,j)\in\mathcal K_n}
|v_{\lambda,ij}v_{\mu,ij}|
\notag\\
&\le
C\sqrt{e_{\lambda,n}e_{\mu,n}},
\label{eq:t3-42}
\end{align}
where
\[
e_{\lambda,n}
=
a_{h_\lambda}^2n^2\rho_nh_\lambda^2.
\]

Consider
\[
Q_n
=
\frac1{2m_n}
\sum_{\lambda\in\Lambda_n}
\sum_{\varsigma\in\{-1,1\}}
P_{\lambda,\varsigma}.
\]
Writing \([x]_+=\max(x,0)\), equations
\eqref{eq:t3-41}--\eqref{eq:t3-42} give
\begin{align}
\chi^2(Q_n,P_0)
&=
\frac1{4m_n^2}
\sum_{\lambda,\mu}
\sum_{\varsigma,\varsigma'}
\left\{
E_0
(L_{\lambda,\varsigma}L_{\mu,\varsigma'})
-1
\right\}
\notag\\
&\le
\frac1{4m_n^2}
\sum_{\substack{\lambda,\mu:\\
\mathcal D_\lambda\cap\mathcal D_\mu\ne\varnothing}}
\sum_{\varsigma,\varsigma'}
\left[
E_0
(L_{\lambda,\varsigma}L_{\mu,\varsigma'})
-1
\right]_+
\notag\\
&\le
\frac{N_{{\rm ov},n}}{m_n^2}
\exp\left\{
C\max_{\lambda\in\Lambda_n}e_{\lambda,n}
\right\}
\notag\\
&\le
\frac{d_n}{m_n}
\exp\left\{
C\max_{\lambda\in\Lambda_n}e_{\lambda,n}
\right\}.
\label{eq:t3-43}
\end{align}
The positive-part step is required because opposite-sign cross terms
can be negative.

Under \eqref{eq:t3-26},
\[
\chi^2(Q_n,P_0)
\le
C(2m_n)^{-1+Cc_{\rm low}+o(1)}
\longrightarrow0
\]
for sufficiently small \(c_{\rm low}>0\).  Hence
\[
\|Q_n-P_0\|_{\rm TV}
\le
\frac12\chi^2(Q_n,P_0)^{1/2}
\longrightarrow0.
\]
For every oracle test \(\varphi\),
\[
P_0(\varphi=1)
+
\sup_{\lambda,\varsigma}
P_{\lambda,\varsigma}(\varphi=0)
\ge
1-\|Q_n-P_0\|_{\rm TV}.
\]
Therefore
\[
R_{n,\circ}^{\rm det,orc}
\ge
1-\|Q_n-P_0\|_{\rm TV}.
\]
The experiment comparison gives
\[
\liminf_{n\to\infty}R_n^{\rm det}\ge1.
\]
The test that never rejects has full-experiment risk one, so
\(R_n^{\rm det}\le1\).  Thus
\[
R_n^{\rm det}\longrightarrow1.
\]

For localization, define
\[
P_{\lambda,+}:=P_{\lambda,+1}.
\]
Let \(V\) be uniform on \(\Lambda_n\), and, conditionally on
\(V=\lambda\), generate the oracle data \(Y\) under \(P_{\lambda,+}\).
If
\[
\overline P_n
=
\frac1{m_n}
\sum_{\lambda\in\Lambda_n}P_{\lambda,+},
\]
then
\begin{align}
I(V;Y)
&=
\frac1{m_n}
\sum_{\lambda\in\Lambda_n}
\operatorname{KL}(P_{\lambda,+},\overline P_n)
\notag\\
&\le
\frac1{m_n}
\sum_{\lambda\in\Lambda_n}
\operatorname{KL}(P_{\lambda,+},P_0)
\notag\\
&\le
C\max_{\lambda\in\Lambda_n}e_{\lambda,n}.
\label{eq:t3-44}
\end{align}
Fano's inequality gives
\begin{equation}
\inf_{\check\lambda}
\sup_{\lambda\in\Lambda_n}
P_{\lambda,+}(\check\lambda\ne\lambda)
\ge
1-
\frac{
C\max_\lambda e_{\lambda,n}+\log2
}{
\log m_n
}.
\label{eq:t3-45}
\end{equation}
Since \(\ell_n/\log m_n\to1\), reducing \(c_{\rm low}\) if necessary
makes the right-hand side at least \(c_F>0\).

Put
\[
r_0
=
\frac14\min\{1,\delta_{\rm pack}\},
\qquad
\mathcal C_\lambda
=
\left\{
D\in\mathcal D_n^{\rm dec}:
d_{\rm loc}(D,\lambda)<r_0
\right\}.
\]
These sets are Borel and pairwise disjoint.  Indeed, common membership
first forces
\[
h_D=h_\lambda=h_\mu.
\]
If \(A=|S_\lambda|=|S_\mu|\), the triangle inequality gives
\[
|S_\lambda\triangle S_\mu|<4r_0A,
\]
whereas \eqref{eq:t3-2} gives
\[
|S_\lambda\triangle S_\mu|
\ge2\delta_{\rm pack}A.
\]
Since \(4r_0\le\delta_{\rm pack}<2\delta_{\rm pack}\), this is
impossible.

Decode \(D\) as the unique \(\lambda\) with
\(D\in\mathcal C_\lambda\), if one exists, and use a fixed
deterministic fallback otherwise.  This decoder is measurable because
\(\Lambda_n\) is finite and the sets \(\mathcal C_\lambda\) are
Borel.  Moreover,
\[
d_{\rm loc}(D,\lambda)
\ge
r_0\mathbf 1\{\check\lambda(D)\ne\lambda\}.
\]
Fano's inequality and the experiment comparison now yield
\[
\liminf_{n\to\infty}R_n^{\rm loc}
\ge
r_0c_F>0.
\]
This proves \eqref{eq:t3-27}.
\end{proof}

\end{document}
